\introheading

When a government forces its public buyers to procure from small firms only, do prices rise because fewer firms bid, or because each firm bids less aggressively? Four decades of research on bid preferences and set-asides---from \citet{marion2007} through \citet{krasnokutskaya2011}, \citet{athey2013}, and \citet{nakabayashi2013}, covering billions of dollars in policy evaluation---have not answered the question directly. The reason is data: existing strategies back out the bidding channel from cross-sectional variation in winning prices alone, a design that conflates entry with aggressiveness by construction. The bid distribution itself has remained a black box.

The answer matters. Set-asides for small and medium-sized enterprises (SMEs) are the most widespread procurement-preference instrument worldwide: the European Union channels a large share of public contract value to SMEs under similar arrangements~\citep{pwc2014,loader2013,loader2015}, OECD and SIGMA surveys document comparable regimes across more than one hundred jurisdictions~\citep{oecd2018,sigma2016,hoekman2020}, and Brazilian federal law mandates exclusive SME tenders for items valued below R\$80{,}000 (US\$22{,}900),\footnote{All USD figures in this paper use a reference exchange rate of R\$3.50 per US dollar, the average of the BCB monthly closing rate over the September 2016--August 2019 sample period. The 2016--2019 annual averages were R\$3.49 (2016), R\$3.19 (2017), R\$3.65 (2018) and R\$3.88 (2019).} covering approximately 97\% of firms representing half of formal employment~\citep{bastos2018, thai2017, bosio2022}. A substantial literature shows that these policies boost firm growth, employment, and supplier-network expansion~\citep{mel2008,freedman2013,szerman2023,cardoza2016,fafchamps2014,mckenzie2015,ferraz2016,cabral2017}. But the design of the fix depends on the channel. If the price premium runs through the extensive margin, expanding the SME roster---through subsidies, registration relief, or capacity-building---can lower the fiscal cost without weakening the preference. If instead it runs through the intensive margin, no roster expansion will help: the only lever is to shrink the \emph{scope} of the rule itself. Policy design hinges on the answer---for the procurement officers in Brazil's 5{,}570 municipalities implementing Law~14{,}133/2021, for the 130 countries with SME-preference regimes in place, and for development economists wrestling with the trade-off between competitive efficiency and distributional objectives in public spending.

Theory points both ways. \citet{bulow1996} show that, across a wide class of auction environments, an additional bidder is worth more to the auctioneer than any sophisticated mechanism---a headline result for the extensive margin. \citet{krasnokutskaya2011} formalize bid preferences in the procurement context and show that both margins are operative. Which dominates quantitatively is an empirical question, and outside the US the evidence for SME set-asides is essentially absent~\citep{nakabayashi2013, szucs2024, decarolis2024}.

This paper shows that the price effect is not a pure headcount effect and that an important within-auction bidding channel---the \emph{intensive margin}---is the most plausible reading of what else is going on. Using phase-2 electronic-auction bid-distribution moments from 3.7~million observations on Sao Paulo's BEC platform, I exploit a March~2018 reinterpretation that extended the state's SME-only regime to medical supplies (Group~65), the only previously exempt product group. The switch was driven by a legalistic isonomy opinion from a separate oversight agency (PGE-SP), plausibly exogenous to Group~65 procurement outcomes. Group~65 is the lone treated cohort and the 76 other groups---subject to SME-only since the 2014 federal mandate---are never-treated controls within the sample window, so identification comes from a clean treated-vs-never-treated comparison, free of the ``forbidden'' implicit comparisons that motivate the modern staggered-DiD literature~\citep{callaway2021,sun2021,goodmanbacon2021,borusyak2024}; the~\citet{kim2019} DiDiR framework delivers the same estimand.

Imposing the SME-only rule on Group~65 raises procurement prices by 10--11\% over 18 months (13--14\% at 6 months, 7--10\% in IPCA-deflated real terms), reduces participating firms by 10--11\%, and cuts valid bids by 7--8\%. Winning firms under open tenders sit 6--14~km further from buyers. The mechanism shows up directly in phase-2 bid moments: the winning bid reaches about six percentage points deeper below the reference under open tenders, an effect that survives conditioning on exactly two, three, or five-plus firms. The~\citet{gelbach2016} decomposition corroborates: entry and winner-composition channels partially offset, netting to about five percent of the price reduction, and 95\% is unexplained by those channels, consistent with a within-auction bidding-behavior story. Pairing a discrete policy cutoff with within-auction bid-distribution data sidesteps the identification problems that have constrained cross-sectional studies of bid preferences~\citep{marion2007, krasnokutskaya2011, athey2013}.

Identification varies by outcome. The log-price effect is the core causal estimate: the price placebo is one-quarter to one-third of the headline (two of three cutoffs marginally significant). The log firms and log bids coefficients have pre-trend contamination that placebos and Sun-Abraham diagnostics do not resolve---they enter as suggestive supporting evidence, not as independently identified. The distance outcome is causally clean (null placebos); the bid-distribution moments driving the mechanism section are within-item and independent of the firms-bids pre-trend issue. Three staggered-DiD estimators (CS2021, Sun-Abraham, Goodman-Bacon) deliver the same sign on all four outcomes. A formal cross-group spillover test bounds redirect bias at less than 1\% of the headline, and RAIS-based firm-size, age, and CNAE controls leave the coefficient invariant. With a single switching cohort, classical permutation inference is underpowered ($p=0.342$); the evidence for causality therefore rests on cross-estimator consistency rather than on that test.

The static fiscal cost of the SME-only rule for Group~65 alone is R\$50--74 million (US\$14--21 million) over the 18-month window (R\$50 million real, R\$74 million nominal), or 7--11\% of the group's procurement value.\footnote{Three fiscal accounting variants appear in the paper and all bracket the same object under different aggregations: R\$50--74~M is the headline range (pooled 18-month coefficient applied to R\$689~M pre-period procurement value, real and nominal); R\$59~M uses quartile-specific coefficients~$\hat\beta_q$ applied to within-quartile procurement value (Table~\ref{tab:counterfactual}); R\$70.7~M aggregates the pharma vs.\ non-pharma split specification (Table~\ref{tab:pharma}). The R\$50--74~M range is the least assumption-heavy.} The number is mechanical---a reduced-form price coefficient times pre-period procurement value---and ignores the completion, entry, and reference-price responses that a structural counterfactual would price. Two findings sharpen the policy implications. First, 92\% of that procurement value is concentrated in the top value quartile, so a mechanical value-threshold exemption recovers 90\% of the accounting cost while preserving the SME preference on 75\% of items by count. SMEs lose access to the top-quartile contracts under this reform---so it is not strictly Pareto-improving, and a full welfare analysis would require modeling the behavioral responses---but even in static terms the reform substantially relaxes the efficiency-equity trade-off that motivates SME-only rules. Second, CMED-regulated pharmaceuticals (class 6531 within Group~65) bear more than twice the distortion of the non-pharma medical supplies: the CMED list-price cap does not shield procurement from the ME/EPP preference, consistent with contract prices lying below the cap and the restriction compressing competition within the residual slack. Two Brazilian regulatory regimes---SME procurement preferences and drug price controls---therefore interact perversely, raising public-health spending above what either instrument alone would deliver.

The paper makes four contributions. \emph{First}, direct bid-distribution evidence on the intensive margin: using phase-2 bid moments, I show the winning bid moves about six percentage points further below the buyer's reference price under open tenders (conditional on firm count) and the runner-up-to-winner gap narrows by 2.7 percentage points, with the effect concentrated in auctions with five or more firms. Prior work has inferred the mechanism from winning-price variation alone~\citep{marion2007, athey2013, krasnokutskaya2011, nakabayashi2013}; 95\% of the price effect here lies outside measured entry and composition channels (net of two partially offsetting channels). \emph{Second}, a sharp counterfactual policy-design result: because 92\% of Group-65 procurement value concentrates in the top value quartile, a mechanical value-threshold exemption---applying the SME-only rule only below about R\$7{,}600 (US\$2{,}180)---recovers 90\% of the accounting fiscal cost while preserving the preference on 75\% of items by count. \emph{Third}, a perverse interaction between two Brazilian price-regulation regimes: CMED-regulated pharmaceuticals (class 6531) bear more than twice the distortion of non-pharma medical supplies, because contract prices sit below the CMED cap and the SME restriction compresses competition within the residual slack---new to the procurement-preference literature~\citep{bosio2022, szucs2024, decarolis2024, colonnelli2022}. \emph{Fourth}, quasi-experimental estimates of SME set-aside costs from a large developing economy (R\$50--74 million, US\$14--21 million, over 18 months for a single product group), rare outside the US and large-firm Europe~\citep{athey2013, krasnokutskaya2011, bandiera2009, best2023}, using a single-cohort \emph{staggered-DiD reversal}~\citep{kim2019, athey2022} against 76 never-treated controls that is robust to Callaway-Sant'Anna, Sun-Abraham, Goodman-Bacon, a cross-group spillover test, and winner-level RAIS controls for firm size, age, and sector.

Section~2 provides the institutional background. Section~3 describes the data. Section~4 presents the empirical analysis, heterogeneity, robustness, and fiscal cost quantification. Section~5 concludes with policy implications.
